% =============================================================================
% Dokumen utama buku ajar (§25, §42).
%
% `main.tex` adalah SATU-SATUNYA berkas LaTeX yang melihat seluruh buku:
% potongan bab di `chapters/chapterNN.tex` tidak berdiri sendiri (tidak ada
% \documentclass di dalamnya) dan tidak mungkin dikompilasi terpisah. Itulah
% sebabnya kompilasi adalah pekerjaan tingkat buku — perintah `export --latex` —
% dan bukan pekerjaan salah satu gate (§26).
%
% Perintah `export --latex` menyalin berkas ini ke `output/latex/` lalu:
%   * mengganti baris \title dan \author di bawah dengan judul buku yang
%     sesungguhnya, dan
%   * mengisi blok \include dengan bab-bab yang sudah mencapai status
%     LATEX_COMPILED, satu baris per bab, berurutan.
%
% Berkas ini dapat ditimpa pengguna lewat `input/source_latex/main.tex`
% (argumen `--latex-template`, §42): prodi yang punya template sendiri hanya
% perlu menyediakan berkas dengan nama yang sama.
% =============================================================================

\documentclass[11pt,a4paper,oneside]{book}
% =============================================================================
% Preamble buku ajar (§25).
%
% Berkas ini dimuat `main.tex` SEBELUM seluruh bab. Ia bukan dokumen: tidak ada
% \documentclass, tidak ada \begin{document}, dan ia tidak boleh di-\include dari
% potongan bab mana pun.
%
% Berkas ini adalah separuh kontrak; separuh lainnya ada di
% `prompts/latex.chapter.md`. Environment yang ditawarkan kepada model di sana
% harus ada di sini, karena environment yang tidak didefinisikan baru ketahuan
% saat kompilasi (§26) — dan pesan galatnya menunjuk ke tempat yang salah, bukan
% ke baris yang menyebabkannya.
%
% Yang harus disunting bersamaan bila salah satu berubah:
%   * nama daftar pustaka  -> `latex/artifacts.py::BIBLIOGRAPHY_FILENAME`
%     (dikunci tes; BibTeX mencari `references.bib` dari nama di
%      \bibliography{references} yang ada di `main.tex`)
%   * \include bab         -> `domain/latex.py::chapter_latex_filename`
% =============================================================================

\usepackage[utf8]{inputenc}
\usepackage[T1]{fontenc}
\usepackage[indonesian]{babel}
\usepackage{amsmath}
\usepackage{amssymb}
\usepackage{amsthm}
\usepackage{graphicx}
\usepackage{booktabs}
\usepackage{enumitem}
\usepackage{xcolor}
\usepackage[colorlinks=true,linkcolor=black,urlcolor=blue]{hyperref}

% listing: satu-satunya environment untuk kode program (§19). Warnanya sengaja
% redup — kode ada di buku untuk dibaca, bukan untuk menarik mata, dan latar
% kelabu pekat membuatnya sulit dicetak.
\usepackage{listings}
\definecolor{codegray}{gray}{0.95}
\lstset{
  basicstyle=\ttfamily\small,
  backgroundcolor=\color{codegray},
  breaklines=true,
  frame=single,
  columns=fullflexible,
  keepspaces=true,
  showstringspaces=false,
}

% definition, example, dan theorem: tiga environment bernama seperti istilah
% blueprint sendiri (§11, §23). Nomornya mengikuti bab, bukan mengikuti
% dokumen, supaya "Definisi 3.2" selalu berarti definisi kedua di bab 3 —
% rujukan silang lintas bab menjadi mungkin dibaca manusia.
\theoremstyle{definition}
\newtheorem{definition}{Definisi}[chapter]
\newtheorem{example}{Contoh}[chapter]
\newtheorem{theorem}{Teorema}[chapter]

% Daftar pustaka. Gaya `plain` dipilih karena ia ada di setiap distribusi
% LaTeX; entri `@misc` yang dihasilkan §25 tidak memuat penulis maupun tahun
% (lihat `domain/latex.py`), dan gaya yang menuntut keduanya akan mencetak
% tanda tanya di tempat yang seharusnya berisi.
\bibliographystyle{plain}


\title{Judul Buku Ajar}
\author{}
\date{\today}

\begin{document}

\frontmatter
\maketitle
\tableofcontents

\mainmatter

% --- bab-bab: diisi `export --latex`, jangan disunting tangan ---------------
% \include{chapters/chapter01}
% \include{chapters/chapter02}
% ---------------------------------------------------------------------------

\backmatter

% Nama `references` di sini yang menentukan nama berkas daftar pustaka:
% BibTeX mencarinya sebagai `references.bib`, dan berkas itulah yang ditulis
% `latex/artifacts.py` (dikunci tes, karena dua tempat ini tidak boleh
% menyimpang tanpa suara).
\bibliography{references}

\end{document}
